\subsection{Units on normal families and on normal valued bases}
\label{subsec:ftb-key-normal-family-units}

We now give the precise extension steps from the field rule and
its DVR units.  The key datum is always the pullback of the
displayed model; no assertion that a generic scalar has the
desired specialization is made before that is proved.

\begin{srlemma}[The source's DVR tests]
\label{lem:ftb-key-dvr-tests}
Let \(A\) be a Noetherian domain with fraction field \(F\).
An element of \(F\) which belongs to every DVR subring of \(F\)
containing \(A\) is integral over \(A\).
Moreover, for primes \(q\subsetneq p\) of a Noetherian ring \(S\),
there are a complete DVR \(O\) and a map \(\alpha:S\to O\) with
\[
                 \ker\alpha=q,\qquad
                 \alpha^{-1}(\mathfrak m_O)=p.
 \tag{U1}\label{eq:ftb-key-dvr-specialization-test}
\]
\end{srlemma}
\begin{proof}
We recall the precise one-dimensional overring calculation
used in the source.
If \(B\) is a one-dimensional Noetherian domain and
\(B\subset C\subset\operatorname{Frac}B\), then \(C\) is
Noetherian of dimension at most one.
For \(0\ne a\in B\), the ring \(B/aB\) is Artinian.
The descending ideals
\[
                    J_n=((a^nC\cap B)+aB)/aB
\]
therefore stabilize.
The resulting equality \(J_{h+1}=J_{h+2}\), for all
sufficiently large \(h\), gives the following downward
step:
\[
 a^{h+1}x=a^{h+2}u+v,\ v\in B
       \quad\Longrightarrow\quad
 a^hx=a^{h+1}u'+v',\ v'\in B.
\]
Indeed replace \(v\), modulo \(aB\), by an element of
\(a^{h+2}C\cap B\), then cancel \(a\).
For each \(x\in C\), clear a denominator \(c\in B\setminus\{0\}\).
Artinianity of \(B/cB\) makes a power of \(a\) divisible
by its next power in that quotient, and gives the first
such displayed expression.
The upward multiplication step and then the downward
steps bring it to one fixed exponent \(n\).
Consequently the \(B\)-module
\[
                        \{x\in C:a^nx\in B\}
\]
surjects onto \(C/aC\).
Multiplication by \(a^n\) identifies this module with a
submodule of \(B\), so it is finite; hence \(C/aC\) is
finite over \(B\).
Every nonzero ideal \(I\subset C\) meets \(B\) nontrivially,
by clearing the denominator of a nonzero element of \(I\).
Choose \(0\ne a\in I\cap B\).
The preceding finite quotient and the principal submodule
\(aC\) show that \(I\) is finitely generated.
Likewise each nonzero prime of \(C\) has a quotient integral
over a finite quotient of \(B\), whose contracting prime
is nonzero and therefore maximal.
That prime of \(C\) is maximal.
This proves both assertions about the overring \(C\).

If \(z\in F\) is not integral over \(A\), the valuation
criterion for integral closure supplies a valuation ring
\(W\subset F\) containing \(A\), with \(z\notin W\).
Then \(y=z^{-1}\) is nonzero and is a nonunit of \(W\).
Put \(B=A[y]\subset W\); its fraction field is \(F\),
and \(y^{-1}\notin B\).
Choose a prime minimal over \(yB\).
The principal ideal theorem makes its localization
one-dimensional and Noetherian.
A valuation overring dominating this localization, with
\(y\) a nonunit, exists by the local valuation-ring
criterion.
The one-dimensional calculation just proved makes this
valuation overring Noetherian.
A Noetherian valuation domain which is not a field is a
DVR.  It contains \(A\), and \(z\) does not belong to it,
which proves the first assertion.

For the second assertion replace \(S\) by
\((S/q)_{p/q}\), a Noetherian local domain of positive
dimension.  Choose a valuation overring \(W\) dominating
it.  Let \(x_1,\ldots,x_r\) generate its maximal ideal.
Choose a nonzero \(y\) among their images having largest
multiplicative valuation, so each \(x_i/y\in W\).
Adjoin these finitely many ratios.
The resulting ring \(B\subset W\) is Noetherian,
birational, and still has \(y\) as a nonunit.
The preceding minimal-prime and overring construction
gives a DVR containing \(B\) in which \(y\) is a nonunit.
All \(x_i=y(x_i/y)\) are then nonunits, so this DVR
dominates the original local ring.
Its completion is again a DVR and the original DVR
injects into it: the intersection of its maximal-ideal
powers is zero.
The completed map is local and injective on
\((S/q)_{p/q}\).
Composing with \(S\to(S/q)_{p/q}\) gives exactly
\eqref{eq:ftb-key-dvr-specialization-test}.
\end{proof}

\begin{srproposition}[The unit and its values at every model point]
\label{prop:ftb-key-family-unit}
Let \(\kappa\) be a field rule with
\(\mathrm{Funct}\) and \(\mathrm{UnitSp}\).
For a normal Noetherian characteristic-zero domain \(S\)
and a key datum \(y/S\), there is a unique \(u\in S^\times\)
such that, for every characteristic-zero field \(K\), map
\(f:S\to K\), and actual key pullback \(x=f^*y\),
\[
                                      \kappa_K(x)=f(u).
 \tag{U2}\label{eq:ftb-key-family-all-points}
\]
If \(\kappa\) also satisfies exponent-one homogeneity and
two such key frames differ by theta scalar \(a\) and
Lie matrix \(M\), their representing values satisfy
\[
                                  a^8u'=\det(M)^4u.
 \tag{U3}\label{eq:ftb-key-family-homog}
\]
\end{srproposition}
\begin{proof}
Let \(F=\operatorname{Frac}S\), and first evaluate the
generic key fibre.
At every height-one prime \(p\), the ring \(S_p\) is a DVR.
Its unit assertion puts both the generic value and its
inverse in \(S_p\).
The height-one intersection identity for the normal
Noetherian ring \(S\) puts both in \(S\).
Their product is \(1\) by injection into \(F\), giving
the unit \(u\).
Functoriality already proves
\eqref{eq:ftb-key-family-all-points} for every injective \(f\).

Now let \(p=\ker f\ne0\).
Use Lemma~\ref{lem:ftb-key-dvr-tests} with \(q=0\), giving
\(\alpha:S\hookrightarrow O\) centered at \(p\).
Pull \(y\) to \(O\).
Its DVR unit \(v\) has the same generic field value as
\(\alpha(u)\), by functoriality; injectivity into
\(\operatorname{Frac}O\) makes \(v=\alpha(u)\) in \(O\).
The residue field of \(O\) has characteristic zero:
the contraction of its maximal ideal is \(p\), and
\(S/p\) embeds in the characteristic-zero field \(K\).
The \(\mathrm{UnitSp}\) specialization therefore gives
the value \(\overline{\alpha(u)}\) in that residue field.

Both the residue map \(S\to O/\mathfrak m_O\) and \(f\)
have kernel \(p\).
They factor through embeddings of
\(\operatorname{Frac}(S/p)\).
Compare the key datum over this fraction field with its
two actual field pullbacks.
Functoriality and injectivity of its residue-field
embedding give value \(u\bmod p\) there, and then value
\(f(u)\) in \(K\).
This proves \eqref{eq:ftb-key-family-all-points} at every
point.  Uniqueness follows at the generic point.
For \eqref{eq:ftb-key-family-homog}, apply the field
homogeneity identity to the same two generic pullbacks,
then inject \(S\) into its fraction field.
\end{proof}

\begin{srproposition}[Normal valued target, with a Noetherian model]
\label{prop:ftb-key-normal-valued-unit}
Let \(S\) be Noetherian, \(V\) a normal domain with fraction
field \(K\) of characteristic zero, \(g:S\to V\), and
\(y/S\) a key datum.
For an actual generic key pullback
\(x=((V\hookrightarrow K)g)^*y\), there is \(u\in V^\times\)
with
\[
                                  \kappa_K(x)=u
                                  \quad\hbox{in }K.
 \tag{U4}\label{eq:ftb-key-valued-unit}
\]
No Noetherian hypothesis on \(V\) is required.
\end{srproposition}
\begin{proof}
Put \(S_0=S/\ker g\) and \(F_0=\operatorname{Frac}S_0\).
These are a Noetherian domain and its fraction field;
\(F_0\) embeds in \(K\).
Pull \(y\) to \(S_0\) and then to \(F_0\), and let \(b\)
be its field value.
For every DVR subring \(W\subset F_0\) containing \(S_0\),
there is the actual key pullback over \(W\).
Its DVR unit shows that \(b,b^{-1}\in W\).
Lemma~\ref{lem:ftb-key-dvr-tests} makes both integral over
\(S_0\).
Functoriality identifies their images in \(K\) with
\(\kappa_K(x)\) and its inverse.
Transport their monic equations through \(S_0\to V\).
Both images are integral over \(V\), and normality of
\(V\) puts them in \(V\).
Their product is \(1\) by injection into \(K\).
They therefore give the asserted unit.
In particular the Noetherian model is used before
passing to the possibly non-Noetherian valued target.
\end{proof}

The exponent-one field rule, its exact norm
\eqref{eq:ftb-key-archnorm}, and these two propositions
supply the complete family interface
\(\mathrm{KeyFam}(g\mapsto1)\), once the previously
retained actual fine-family and analytic-realization
inputs have been joined.
The fields, DVRs, normal Noetherian families, and normal
valued targets in that interface have exactly the scopes
printed here.

\paragraph{Frozen source concordance.}
The DVR approximation, one-dimensional overring and
completion mechanisms are
\srcid{src/Development/AbelianVarietyTheory/Proofs/FoundationProofs/%
RangeAlgebraMapEqIInfDiscreteValuationSubring.lean}
(\srcid{d373da2375cb8f30eae09b8bab0170253dbd79a0}).
The explicit integrality test is
\srcid{IsIntegralOfForallMemDiscreteValuationSubring.lean}
(\srcid{37bf7d4a155cf39f6b2a69dba9a64553e981a429})
in that directory.
The remaining prefix is
\srcid{src/Development/AbelianVarietyTheory/Proofs/ThetaLawProofs/}.
The height-one generic-unit construction is
\srcid{ExistsUnitThetaKeyRuleSmoothBase.lean}
(\srcid{36413650a428826ff71462b79aded750faaa2d18}),
and the every-point DVR test is
\srcid{IsValueOfGeneric.lean}
(\srcid{f3f51dfb486c4f00e04e1c527e9d5264b4e386e8}).
The exact assemblies are
\srcid{FamUnit.lean}
(\srcid{fea31937eccfed5f0483942977d688d14d9dc59d}),
\srcid{FamHomogOfHomog.lean}
(\srcid{a28c119b27b4f1d6cf94d2c40f04e7f9fb4ef103}),
\srcid{ValueMemValuationSubring.lean}
(\srcid{ac1011e3619e36a8541b57dc4c001a258576baeb}),
\srcid{ThetaKeyRuleValuationUnitOfFunctorial.lean}
(\srcid{97457f92dcd20abfde83b5eafe6c49c7940481c9}),
and \srcid{KeyFamOfKeyFormula.lean}
(\srcid{660b0c8e7c854d703db866755e2f6b4cc423274b}).
These arguments add no new moduli or auxiliary analytic
hypothesis to the field rule; the field rule's own
separate producers remain as stated earlier.
